\documentclass[12pt,a4paper]{article}

\usepackage[a4paper,margin=18mm]{geometry}
\usepackage[T2A]{fontenc}
\usepackage[utf8]{inputenc}
\usepackage[russian]{babel}
\usepackage{amsmath,amssymb,mathtools}
\usepackage{array,booktabs,longtable}
\usepackage{xcolor}
\usepackage{hyperref}
\usepackage{tikz}

\hypersetup{
  colorlinks=true,
  linkcolor=blue!55!black,
  urlcolor=blue!55!black
}

\setlength{\parindent}{0pt}
\setlength{\parskip}{5pt}
\renewcommand{\arraystretch}{1.25}

\newcommand{\reviewhead}[1]{%
  \par\medskip
  \noindent\textbf{#1}\par\smallskip
}

\begin{document}

\begin{titlepage}
\thispagestyle{empty}
\begin{center}
\vspace*{0.24\textheight}

{\LARGE\bfseries Ревью домашней работы\par}

\vspace{18mm}

{\large Автор: Лёвин Артём Александрович\par}

\vspace{7mm}

{\large 9 сентября 2026 года\par}

\vfill

\begin{tikzpicture}
\draw[blue!35!black,line width=0.7pt]
(-3.5,0) .. controls (-2.0,0.8) and (-0.8,-0.5) .. (0,0)
        .. controls (0.8,0.5) and (2.0,-0.8) .. (3.5,0);
\end{tikzpicture}

\vspace{7mm}
\end{center}
\end{titlepage}

\newpage
\section*{Сводная таблица}
\small
\begin{center}
\begin{tabular}{@{}p{9mm}p{30mm}p{45mm}p{34mm}p{34mm}@{}}
\toprule
\textbf{№} &
\textbf{Тема} &
\textbf{Суть ошибки} &
\textbf{Ваш ответ} &
\textbf{Правильный ответ} \\
\midrule
\hyperref[task:3]{3} &
Линейная функция: коэффициент и пересечения с осями &
Перепутаны точки пересечения с осями \(Oy\) и \(Ox\). &
\(k=-3\); функция убывает; для \(Oy\): \((2;0)\); для \(Ox\): \((0;6)\). &
\(k<0\); функция убывает; \(Oy:(0;6)\); \(Ox:(2;0)\). \\
\addlinespace
\hyperref[task:5]{5} &
Уравнение прямой по двум точкам &
Решение начато, но коэффициенты \(k,b\), уравнение и монотонность не найдены. &
Не указан. &
\(y=-x+3\); функция убывает. \\
\addlinespace
\hyperref[task:6]{6} &
Параметр в линейной функции &
В пункте 1 анализируется знак \(m\), хотя знак наклона задаёт \(k=m-2\). &
1) \(m<0\); \quad 2) \(m=4\). &
1) \(m>2\); \quad 2) \(m=4\). \\
\addlinespace
\hyperref[task:10]{10} &
Линейная функция, пересечения с осями, неравенство &
Решение и ответ в приложенных материалах не представлены. &
Не представлен. &
\(y=-x+4\); \(Oy:(0;4)\); \(Ox:(4;0)\); \(x\le 2\), то есть \((-\infty;2]\). \\
\bottomrule
\end{tabular}
\end{center}
\normalsize

\newpage
\vspace*{\fill}
\begin{center}
{\LARGE\bfseries Разбор ошибок}
\end{center}
\vspace*{\fill}
\newpage

\phantomsection\label{task:3}
\section*{Задание 3}

\reviewhead{Текст задания}
Дана функция
\[
y=-3x+6.
\]
Определите:
\begin{enumerate}
\item знак \(k\);
\item возрастает или убывает функция;
\item точку пересечения с осью \(Oy\);
\item точку пересечения с осью \(Ox\).
\end{enumerate}

\reviewhead{Ваше решение}
В работе записано:
\[
k=-3,
\]
далее указано, что функция убывает, так как \(k<0\). Затем для третьего пункта записана точка
\[
(2;0),
\]
а для четвёртого пункта ---
\[
(0;6).
\]

\reviewhead{Ваш ответ}
\[
k=-3,\qquad \text{функция убывает},\qquad
Oy:(2;0),\qquad Ox:(0;6).
\]

\reviewhead{Ошибка}
\textcolor{red}{Ошибка:} точки пересечения с осями \(Oy\) и \(Ox\) записаны в обратном порядке.

\reviewhead{Где именно возникает ошибка}
До определения монотонности решение идёт верно:
\[
k=-3<0,
\]
поэтому функция действительно убывает.

Первый неверный шаг появляется в пункте о пересечении с осью \(Oy\):
\[
Oy:(2;0).
\]
Эта точка не может лежать на оси \(Oy\), потому что у любой точки оси \(Oy\) первая координата обязательно равна нулю.

\reviewhead{Почему этот шаг неверен}
Для пересечения с осью \(Oy\) нужно положить
\[
x=0,
\]
поскольку ось \(Oy\) состоит из всех точек вида \((0;y)\).

Для пересечения с осью \(Ox\), наоборот, нужно положить
\[
y=0,
\]
поскольку ось \(Ox\) состоит из всех точек вида \((x;0)\).

Точка \((2;0)\) имеет нулевую вторую координату, поэтому она лежит на \(Ox\), а не на \(Oy\).
Точка \((0;6)\) имеет нулевую первую координату, поэтому она лежит на \(Oy\), а не на \(Ox\).

\reviewhead{Ключевая идея}
\textcolor{blue!60!black}{Ключевая идея:}
для точки на \(Oy\) всегда \(x=0\); для точки на \(Ox\) всегда \(y=0\).

\reviewhead{Исправление}
Начинаем с последнего верного места.

Коэффициент при \(x\):
\[
k=-3<0,
\]
следовательно, функция убывает.

Пересечение с \(Oy\):
\[
x=0,
\]
\[
y=-3\cdot 0+6=6.
\]
Поэтому
\[
Oy:(0;6).
\]

Пересечение с \(Ox\):
\[
y=0,
\]
\[
-3x+6=0,
\]
\[
-3x=-6,
\]
\[
x=2.
\]
Поэтому
\[
Ox:(2;0).
\]

\reviewhead{Полное правильное решение}
Функция имеет вид
\[
y=kx+b,
\]
где
\[
k=-3,\qquad b=6.
\]

Так как
\[
k=-3<0,
\]
линейная функция убывает.

Для пересечения с осью \(Oy\) полагаем \(x=0\):
\[
y=-3\cdot 0+6=6.
\]
Получаем точку
\[
(0;6).
\]

Для пересечения с осью \(Ox\) полагаем \(y=0\):
\[
0=-3x+6,
\]
\[
3x=6,
\]
\[
x=2.
\]
Получаем точку
\[
(2;0).
\]

\reviewhead{Проверка}
Подставим обе найденные точки в исходную формулу.

Для \((0;6)\):
\[
6=-3\cdot 0+6=6.
\]

Для \((2;0)\):
\[
0=-3\cdot 2+6=-6+6=0.
\]

Обе точки принадлежат графику, и каждая лежит на требуемой оси.

\reviewhead{Итог}
\textcolor{teal!60!black}{Итог:}
\[
k<0,\qquad \text{функция убывает},\qquad
Oy:(0;6),\qquad Ox:(2;0).
\]

\reviewhead{Задача на закрепление}
Дана функция
\[
y=2x-8.
\]
Определите знак \(k\), укажите, возрастает или убывает функция, и найдите точки её пересечения с осями \(Oy\) и \(Ox\).

\newpage
\phantomsection\label{task:5}
\section*{Задание 5}

\reviewhead{Текст задания}
Прямая проходит через точки
\[
C(-2;5),\qquad D(4;-1).
\]
Найдите её уравнение и укажите, возрастает или убывает соответствующая функция.

\reviewhead{Ваше решение}
В работе выписаны точки
\[
C(-2;5),\qquad D(4;-1)
\]
и общий вид линейной функции
\[
y=kx+b.
\]
Дальнейшие вычисления для определения \(k\) и \(b\) не представлены.

\reviewhead{Ваш ответ}
Не указан.

\reviewhead{Ошибка}
\textcolor{red}{Ошибка:} решение не доведено до нахождения конкретного уравнения прямой и определения её монотонности.

\reviewhead{Где именно возникает ошибка}
Последний корректный шаг:
\[
y=kx+b.
\]
Это правильный общий вид линейной функции.

Первый неполный шаг состоит в том, что после этой записи координаты точек \(C\) и \(D\) не подставлены в формулу. Поэтому коэффициенты \(k\) и \(b\) остаются неизвестными, а требуемое уравнение не получено.

\reviewhead{Почему это неверно}
Запись
\[
y=kx+b
\]
задаёт не одну конкретную прямую, а целое семейство прямых.

Чтобы выбрать именно ту прямую, которая проходит через обе заданные точки, необходимо использовать оба условия:
\[
C(-2;5)\in y=kx+b,
\]
\[
D(4;-1)\in y=kx+b.
\]
Каждая точка даёт одно уравнение для \(k\) и \(b\). Два уравнения позволяют найти оба неизвестных коэффициента.

После нахождения \(k\) знак этого коэффициента определяет монотонность линейной функции:
\[
k>0 \Rightarrow \text{функция возрастает},
\]
\[
k<0 \Rightarrow \text{функция убывает}.
\]

\reviewhead{Ключевая идея}
\textcolor{blue!60!black}{Ключевая идея:}
две заданные точки дают две подстановки в \(y=kx+b\); из полученной системы находятся \(k\) и \(b\).

\reviewhead{Исправление}
Продолжаем с уже записанного верного общего вида:
\[
y=kx+b.
\]

Подставляем \(C(-2;5)\):
\[
5=-2k+b.
\]

Подставляем \(D(4;-1)\):
\[
-1=4k+b.
\]

Получаем систему
\[
\begin{cases}
5=-2k+b,\\
-1=4k+b.
\end{cases}
\]

Вычтем второе уравнение из первого:
\[
5-(-1)=(-2k+b)-(4k+b),
\]
\[
6=-6k,
\]
\[
k=-1.
\]

Подставим \(k=-1\) в первое уравнение:
\[
5=-2(-1)+b,
\]
\[
5=2+b,
\]
\[
b=3.
\]

Следовательно,
\[
y=-x+3.
\]

Так как
\[
k=-1<0,
\]
функция убывает.

\reviewhead{Полное правильное решение}
Пусть уравнение прямой имеет вид
\[
y=kx+b.
\]

Так как прямая проходит через \(C(-2;5)\), получаем:
\[
5=-2k+b.
\]

Так как прямая проходит через \(D(4;-1)\), получаем:
\[
-1=4k+b.
\]

Решаем систему:
\[
\begin{cases}
5=-2k+b,\\
-1=4k+b.
\end{cases}
\]

Из первого уравнения вычтем второе:
\[
6=-6k,
\]
\[
k=-1.
\]

Теперь найдём \(b\):
\[
5=-2(-1)+b,
\]
\[
5=2+b,
\]
\[
b=3.
\]

Поэтому искомое уравнение:
\[
\boxed{y=-x+3}.
\]

Коэффициент при \(x\) равен
\[
k=-1<0,
\]
значит, соответствующая линейная функция убывает.

\reviewhead{Проверка}
Проверим обе заданные точки.

Для \(C(-2;5)\):
\[
-(-2)+3=2+3=5.
\]

Для \(D(4;-1)\):
\[
-4+3=-1.
\]

Обе точки удовлетворяют уравнению \(y=-x+3\).

\reviewhead{Итог}
\textcolor{teal!60!black}{Итог:}
\[
y=-x+3,\qquad \text{функция убывает}.
\]

\reviewhead{Задача на закрепление}
Прямая проходит через точки
\[
P(-1;4),\qquad Q(3;-4).
\]
Найдите её уравнение и укажите, возрастает или убывает соответствующая функция.

\newpage
\phantomsection\label{task:6}
\section*{Задание 6}

\reviewhead{Текст задания}
Дана функция
\[
y=(m-2)x+5.
\]
\begin{enumerate}
\item При каких \(m\) функция возрастает?
\item Найдите \(m\), если график проходит через точку \(P(3;11)\).
\end{enumerate}

\reviewhead{Ваше решение}
В первом пункте записано:
\[
m<0.
\]

Во втором пункте выполнена подстановка координат точки \(P(3;11)\):
\[
11=(m-2)\cdot 3+5,
\]
далее получено
\[
11=3m-6+5,
\]
\[
3m=12,
\]
\[
m=4.
\]

\reviewhead{Ваш ответ}
\[
1)\ m<0;\qquad 2)\ m=4.
\]

\reviewhead{Ошибка}
\textcolor{red}{Ошибка:} в первом пункте рассматривается знак параметра \(m\), а нужно рассматривать знак коэффициента при \(x\), то есть \(m-2\).

\reviewhead{Где именно возникает ошибка}
Из формулы
\[
y=(m-2)x+5
\]
видно, что коэффициент наклона равен
\[
k=m-2.
\]

Первый неверный шаг:
\[
m<0.
\]
Это условие не связано напрямую с возрастанием функции. Например, при \(m=1\) имеем
\[
k=1-2=-1<0,
\]
и функция убывает. Значит, проверять знак самого \(m\) недостаточно.

Второй пункт решён верно и может быть сохранён.

\reviewhead{Почему этот шаг неверен}
Для линейной функции
\[
y=kx+b
\]
монотонность определяется только знаком коэффициента \(k\):
\[
k>0 \Rightarrow \text{функция возрастает},
\]
\[
k<0 \Rightarrow \text{функция убывает}.
\]

В данном задании
\[
k=m-2,
\]
поэтому условие возрастания имеет вид
\[
m-2>0,
\]
а не \(m<0\).

\reviewhead{Ключевая идея}
\textcolor{blue!60!black}{Ключевая идея:}
сначала нужно выделить коэффициент при \(x\), а уже затем исследовать его знак.

\reviewhead{Исправление}
Для первого пункта:
\[
k=m-2.
\]
Функция возрастает, если
\[
k>0.
\]
Значит,
\[
m-2>0,
\]
\[
m>2.
\]

Второй пункт оставляем, так как вычисления выполнены правильно:
\[
11=(m-2)\cdot 3+5,
\]
\[
11=3m-6+5,
\]
\[
11=3m-1,
\]
\[
3m=12,
\]
\[
m=4.
\]

\reviewhead{Полное правильное решение}
\textbf{1) Условие возрастания.}

Функция
\[
y=(m-2)x+5
\]
линейная, и её коэффициент при \(x\):
\[
k=m-2.
\]

Линейная функция возрастает тогда и только тогда, когда
\[
k>0.
\]

Поэтому:
\[
m-2>0,
\]
\[
\boxed{m>2}.
\]

\textbf{2) Условие прохождения через \(P(3;11)\).}

Если точка \(P(3;11)\) лежит на графике, её координаты удовлетворяют уравнению:
\[
11=(m-2)\cdot 3+5.
\]

Раскрываем скобки:
\[
11=3m-6+5,
\]
\[
11=3m-1.
\]

Переносим \(-1\) в левую часть:
\[
12=3m,
\]
\[
\boxed{m=4}.
\]

\reviewhead{Проверка}
При \(m=4\):
\[
y=(4-2)x+5=2x+5.
\]

Для \(x=3\):
\[
y=2\cdot 3+5=11,
\]
поэтому точка \(P(3;11)\) действительно лежит на графике.

Кроме того,
\[
4>2,
\]
поэтому при найденном во втором пункте значении \(m=4\) функция действительно возрастает.

\reviewhead{Итог}
\textcolor{teal!60!black}{Итог:}
\[
1)\ m>2;\qquad 2)\ m=4.
\]

\reviewhead{Задача на закрепление}
Дана функция
\[
y=(a+3)x-4.
\]
Определите, при каких \(a\) функция возрастает, и найдите \(a\), если её график проходит через точку \(Q(2;6)\).

\newpage
\phantomsection\label{task:10}
\section*{Задание 10}

\reviewhead{Текст задания}
Прямая проходит через точки
\[
A(-3;7),\qquad B(5;-1).
\]
\begin{enumerate}
\item Найдите уравнение прямой.
\item Найдите точки её пересечения с осями координат.
\item Решите неравенство
\[
y\ge 2,
\]
если \(y\) задаётся найденной линейной функцией.
\item Запишите множество подходящих значений \(x\) промежутком.
\end{enumerate}

\reviewhead{Ваше решение}
В приложенных материалах решение задания 10 не представлено.

\reviewhead{Ваш ответ}
Не представлен.

\reviewhead{Ошибка}
\textcolor{red}{Ошибка:} обязательное решение задания отсутствует в представленных материалах.

\reviewhead{Где именно возникает ошибка}
Последний корректный шаг отсутствует, поскольку решение задания не начато в доступных материалах.

Первый неполный шаг --- отсутствие определения коэффициентов прямой по двум заданным точкам. Без уравнения прямой невозможно корректно выполнить последующие пункты о пересечениях с осями и неравенстве.

\reviewhead{Почему это неверно}
Задание состоит из связанных пунктов. Результат первого пункта,
то есть конкретная формула линейной функции, используется дальше:
по ней находятся точки пересечения с осями и решается неравенство.

Если уравнение прямой не найдено, дальнейшие пункты не имеют вычислительной основы.

\reviewhead{Ключевая идея}
\textcolor{blue!60!black}{Ключевая идея:}
сначала по двум точкам находим \(k\) и \(b\) в формуле \(y=kx+b\); затем используем найденную функцию во всех остальных пунктах.

\reviewhead{Исправление}
Найдём коэффициент \(k\) по двум точкам:
\[
k=\frac{y_B-y_A}{x_B-x_A}.
\]

Подставляем координаты:
\[
k=\frac{-1-7}{5-(-3)}
=\frac{-8}{8}
=-1.
\]

Теперь найдём \(b\). Подставим точку \(A(-3;7)\) в
\[
y=-x+b:
\]
\[
7=-(-3)+b,
\]
\[
7=3+b,
\]
\[
b=4.
\]

Следовательно,
\[
y=-x+4.
\]

Пересечение с \(Oy\):
\[
x=0,
\]
\[
y=4,
\]
поэтому
\[
(0;4).
\]

Пересечение с \(Ox\):
\[
y=0,
\]
\[
-x+4=0,
\]
\[
x=4,
\]
поэтому
\[
(4;0).
\]

Решаем неравенство:
\[
-x+4\ge 2,
\]
\[
-x\ge -2.
\]
При умножении обеих частей на \(-1\) знак неравенства меняется:
\[
x\le 2.
\]

Промежуток:
\[
(-\infty;2].
\]

\reviewhead{Полное правильное решение}
\textbf{1) Уравнение прямой.}

Пусть
\[
y=kx+b.
\]

Коэффициент наклона:
\[
k=\frac{-1-7}{5-(-3)}
=\frac{-8}{8}
=-1.
\]

Получаем
\[
y=-x+b.
\]

Используем точку \(A(-3;7)\):
\[
7=-(-3)+b,
\]
\[
7=3+b,
\]
\[
b=4.
\]

Значит,
\[
\boxed{y=-x+4}.
\]

\textbf{2) Пересечения с осями.}

С осью \(Oy\): полагаем \(x=0\):
\[
y=-0+4=4.
\]
Точка:
\[
\boxed{(0;4)}.
\]

С осью \(Ox\): полагаем \(y=0\):
\[
0=-x+4,
\]
\[
x=4.
\]
Точка:
\[
\boxed{(4;0)}.
\]

\textbf{3) Решение неравенства \(y\ge 2\).}

Подставляем
\[
y=-x+4:
\]
\[
-x+4\ge 2.
\]

Вычитаем \(4\) из обеих частей:
\[
-x\ge -2.
\]

Умножаем на \(-1\) и обязательно меняем знак неравенства:
\[
\boxed{x\le 2}.
\]

\textbf{4) Запись промежутком.}

Все подходящие значения:
\[
\boxed{(-\infty;2]}.
\]

\reviewhead{Проверка}
Проверим, что найденная прямая проходит через обе исходные точки.

Для \(A(-3;7)\):
\[
-(-3)+4=3+4=7.
\]

Для \(B(5;-1)\):
\[
-5+4=-1.
\]

Проверим граничное значение неравенства \(x=2\):
\[
y=-2+4=2,
\]
поэтому \(x=2\) должно входить в ответ, что согласуется с квадратной скобкой в
\[
(-\infty;2].
\]

Если взять, например, \(x=3\), то
\[
y=-3+4=1<2,
\]
поэтому значения \(x>2\) действительно не подходят.

\reviewhead{Итог}
\textcolor{teal!60!black}{Итог:}
\[
y=-x+4,
\]
\[
Oy:(0;4),\qquad Ox:(4;0),
\]
\[
x\le 2,\qquad (-\infty;2].
\]

\reviewhead{Задача на закрепление}
Прямая проходит через точки
\[
M(-1;5),\qquad N(3;-3).
\]
Найдите её уравнение, найдите точки пересечения с осями координат, решите неравенство \(y\ge 1\) для найденной линейной функции и запишите множество подходящих значений \(x\) промежутком.

\end{document}
